\documentclass[11pt,a4paper]{article}

% ============================================================
% PE4_U4 - Práctica Experimental Unidad IV
% Ingeniería de Requerimientos [20303] - ISR-401
% UTEQ - Facultad de Ciencias de la Computación
% Sistema: SGCV-IA
% ============================================================

\usepackage[utf8]{inputenc}
\usepackage[T1]{fontenc}
\usepackage{helvet}
\renewcommand{\familydefault}{\sfdefault}
\usepackage[a4paper,margin=2.2cm,headheight=28pt]{geometry}
\usepackage{xcolor}
\usepackage{fancyhdr}
\usepackage{longtable}
\usepackage{booktabs}
\usepackage{array}
\usepackage{graphicx}
\usepackage{hyperref}
\usepackage{enumitem}
\usepackage{titlesec}
\usepackage{tabularx}
\usepackage{amsmath}

% ---------- Paleta institucional ----------
\definecolor{uteqverde}{HTML}{2C5F2D}
\definecolor{uteqgris}{HTML}{4D4D4D}

% ---------- Encabezado / pie de página ----------
\pagestyle{fancy}
\fancyhf{}
\renewcommand{\headrulewidth}{0.6pt}
\renewcommand{\footrulewidth}{0.4pt}
\fancyhead[L]{\small\color{uteqgris} Ingeniería de Requerimientos (ISR-401) -- UTEQ}
\fancyhead[R]{\small\color{uteqgris} PE4\_U4 -- Unidad IV}
\fancyfoot[C]{\small\color{uteqgris} \thepage}
\fancyfoot[R]{\small\color{uteqgris} SGCV-IA}

% ---------- Comando de sección institucional ----------
% Uso: \seccion{N.}{Título de la sección}
\newcommand{\seccion}[2]{%
  \vspace{0.6cm}
  \noindent
  {\color{white}\colorbox{uteqverde}{%
    \parbox[c][1.05cm]{\dimexpr\textwidth-16pt\relax}{%
      \hspace{6pt}\textbf{\large #1\ #2}%
    }%
  }}%
  \vspace{0.35cm}
  \par
}

\newcommand{\subseccion}[1]{%
  \vspace{0.35cm}
  \noindent{\color{uteqverde}\textbf{\large #1}}\par
  \vspace{0.15cm}
}

% Título de tabla/figura en verde, cuerpo en blanco y negro (sin tcolorbox)
\newcommand{\tituloTabla}[1]{{\color{uteqverde}\textbf{#1}}}

\setlist[itemize]{leftmargin=1.4em}
\setlist[enumerate]{leftmargin=1.6em}

\hypersetup{
  hidelinks,
  pdftitle={PE4\_U4 - Práctica Experimental Unidad IV - SGCV-IA},
  pdfauthor={Equipo SGCV-IA},
  pdfsubject={Ingeniería de Requerimientos ISR-401}
}

\begin{document}

% ================================================================
% 1. CARÁTULA
% ================================================================
\begin{titlepage}
\centering
\vspace*{0.5cm}

% [COMPLETAR: logo institucional opcional -- \includegraphics[width=3.5cm]{figuras/logo_uteq.png}]

{\Large\bfseries UNIVERSIDAD TÉCNICA ESTATAL DE QUEVEDO\par}
\vspace{0.15cm}
{\large Facultad de Ciencias de la Computación\par}
{\large Carrera de Software (Rediseño)\par}

\vspace{1.2cm}
\rule{0.85\textwidth}{1pt}
\vspace{0.5cm}

{\Large\bfseries\color{uteqverde} PE4\_U4 -- Práctica Experimental Unidad IV\par}
\vspace{0.3cm}
{\large Validación, Gestión de Requisitos y Herramientas CASE\par}
\vspace{0.2cm}
{\normalsize\itshape Inspección Fagan · Change Control Board · Trazabilidad en \\
herramienta CASE · Baseline con Git\par}

\vspace{0.5cm}
\rule{0.85\textwidth}{1pt}
\vspace{1cm}

{\large\bfseries Sistema del PFC: SGCV-IA\par}
{\normalsize Sistema de Gestión de Clínica Veterinaria con Inteligencia Artificial\par}

\vspace{1.2cm}

\begin{tabular}{ll}
\textbf{Asignatura:} & Ingeniería de Requerimientos [20303] -- 4.\textordmasculine\ nivel \\
\textbf{Paralelo:} & A \\
\textbf{Período:} & 2026--2027 PPA \\
\textbf{Docente:} & Ing. Gleiston Cicerón Guerrero Ulloa, PhD. \\
\textbf{Semana de entrega:} & 14 \\
\end{tabular}

\vspace{1cm}

\begin{tabular}{ll}
\multicolumn{2}{l}{\textbf{Integrantes y rol desempeñado en la práctica:}}\\[4pt]
Marcillo Ponce, Alberto Jeanpool & Moderador / Inspector 2 \\
Herrera Ramos, Thais Melanie & Inspector 1 \\
Arboleda Yanza, Francisco Javier & Lector \\
\end{tabular}

\vspace{1.2cm}

\textbf{Fecha:} 05/08/2026

\vspace{0.6cm}

% Requisito G1 (gatekeeper): la URL debe quedar visible en UNA SOLA LÍNEA
% [PENDIENTE DEL EQUIPO: reemplazar por la URL pública real del repositorio antes de compilar/entregar -- criterio de piso G1]
\textbf{Repositorio:} \url{https://github.com/[NombreEquipo]/ISR401-PFC-ERS-SGCVIA}

\vfill
{\small Guía y Rúbrica PE4 -- Práctica Experimental Unidad IV -- ISR-401}
\end{titlepage}

% ================================================================
% 2. RESUMEN Y OBJETIVOS
% ================================================================
\seccion{2.}{RESUMEN Y OBJETIVOS}

\subseccion{Resumen}
% Máximo 200 palabras.
Este informe documenta la Práctica Experimental Unidad IV sobre el ERS vigente
(Entrega 2A) del Sistema de Gestión de Clínica Veterinaria con Inteligencia
Artificial (SGCV-IA). El equipo ejecutó una inspección formal tipo Fagan sobre
las secciones 3.2 (RF-01 a RF-25) y 3.3 (RNF-01 a RNF-15) del ERS, con los
cuatro roles diferenciados repartidos entre tres integrantes, preparación
individual documentada en el Anexo A y un registro consolidado de 18 defectos
únicos (0 críticos, 6 mayores, 12 menores; densidad de 1,38 defectos por
página sobre 13 páginas efectivas). El 100\,\% de los defectos mayores fue
corregido con texto antes/después. A partir de uno de esos defectos y de dos
escenarios adicionales, se tramitaron tres RFC de naturaleza distinta (calidad,
alcance y restricción normativa) ante un Change Control Board simulado, que
aprobó los tres --uno sin condiciones y dos con condiciones-- y produjo la
versión v1.1 del ERS. El informe cierra con la comparación de cuatro
herramientas CASE de gestión de requisitos, el contraste entre IR tradicional
e IR ágil aplicado al SGCV-IA, y las conclusiones críticas del equipo. La
trazabilidad en herramienta CASE y la línea base con Git (Pasos 3 y 4 de la
guía) quedan documentadas en las secciones 7 y 8 conforme a la evidencia
disponible al cierre de este documento.

\subseccion{Objetivo general}
Aplicar técnicas formales de validación de requisitos, gestionar cambios mediante un
Change Control Board simulado, implementar la trazabilidad del ERS en una herramienta
CASE, establecer una línea base con Git y comparar metodológicamente la Ingeniería de
Requisitos tradicional frente a la ágil, sobre el ERS real del Sistema de Gestión de
Clínica Veterinaria con Inteligencia Artificial (SGCV-IA).

\subseccion{Objetivos específicos}
\begin{enumerate}
  \item Ejecutar una inspección formal tipo Fagan sobre el ERS del SGCV-IA, con roles
  diferenciados, lista de verificación y registro de defectos.
  \item Calcular e interpretar métricas de inspección (densidad de defectos,
  distribución por tipo y severidad, tasa de corrección).
  \item Tramitar tres solicitudes formales de cambio (RFC) con análisis de impacto
  sustentado en la matriz de trazabilidad.
  \item Deliberar y decidir en un CCB simulado, dejando acta motivada y versión
  actualizada del ERS.
  \item Construir la trazabilidad bidireccional del ERS en una herramienta CASE de
  gestión, con campos personalizados y enlaces verificables.
  \item Establecer y publicar la línea base del ERS con control de versiones.
  \item Comparar herramientas CASE de IR y contrastar IR tradicional frente a IR ágil,
  cerrando con conclusiones críticas propias.
\end{enumerate}

% ================================================================
% 3. FUNDAMENTO TEÓRICO
% ================================================================
\seccion{3.}{FUNDAMENTO TEÓRICO}

\subseccion{3.1 Marco normativo}

La ISO/IEC/IEEE 29148:2018 organiza la ingeniería de requisitos en tres pilares:
los procesos de requisitos (cláusula 6), los ítems de información como el ERS
(cláusula 7) y los criterios de calidad de un requisito individual y del conjunto
(cláusula 5) \cite{ieee29148}. Dentro de la cláusula 6, el subapartado 6.4 ubica la
validación de requisitos como la actividad que confirma que las especificaciones
reflejan fielmente las necesidades del stakeholder, distinguiéndola de la
verificación de la conformidad interna del documento. Esta ubicación normativa es
la que sustenta que la inspección Fagan aplicada en esta práctica se considere una
técnica de validación y no de simple corrección de forma.

El control de cambios sobre los requisitos, por su parte, se apoya en dos normas
complementarias a la 29148. La ISO/IEC/IEEE 15288:2023 define el proceso de gestión
de la configuración a nivel de sistema, del cual se deriva la necesidad de una
línea base auditable y de un registro de decisiones sobre cada solicitud de cambio
\cite{ieee15288}. La ISO/IEC/IEEE 12207:2017 traslada ese mismo proceso al ámbito
específico del ciclo de vida del software, detallando las actividades de control de
cambios y de gestión de la configuración que un equipo de desarrollo debe ejecutar
sobre sus artefactos, entre ellos el ERS \cite{ieee12207}.

Para evaluar la calidad del ERS resultante, esta práctica aplica el modelo de
calidad de producto de la ISO/IEC 25010:2023, que reemplaza al de 2011 y organiza
la calidad en nueve características: adecuación funcional, eficiencia de
desempeño, compatibilidad, capacidad de interacción, fiabilidad, seguridad,
mantenibilidad, flexibilidad y seguridad de uso (\textit{safety}), esta última
incorporada como característica nueva en la versión 2023 \cite{iso25010}. De estas
nueve, la mantenibilidad y la capacidad de interacción son las más relevantes al
momento de clasificar los defectos de ambigüedad y de verificabilidad detectados
en la inspección, mientras que la seguridad y la fiabilidad enmarcan los NFR que
se modifican en las solicitudes de cambio del Paso 2.

Complementan este marco el SWEBOK v4.0, cuya área de conocimiento de requisitos de
software (capítulo 1) dedica sus apartados 1.6 a 1.8 a la validación, la gestión
del cambio y la trazabilidad de requisitos como actividades transversales al resto
del proceso \cite{swebok2024}; la Parte V de Pohl (2025), centrada en la gestión y
evolución de requisitos a lo largo del ciclo de vida \cite{pohl2025}; y la
séptima edición del PMBOK, que define al Change Control Board (CCB) como el
órgano formal responsable de evaluar, aprobar, rechazar o diferir las solicitudes
de cambio sobre los artefactos de un proyecto, rol que en esta práctica se simula
sobre el ERS del SGCV-IA \cite{pmbok2021}.

\subseccion{3.2 Validación de requisitos}

La verificación y la validación son actividades distintas y complementarias.
La verificación responde a la pregunta \textit{¿el ERS está construido
correctamente?}, es decir, si cumple con el formato, la trazabilidad interna y los
criterios de calidad de la 29148; la validación responde a \textit{¿se está
construyendo el ERS correcto?}, esto es, si los requisitos documentados
representan realmente las necesidades del stakeholder \cite{ieee29148}. La
inspección Fagan, técnica central de esta práctica, cumple ambas funciones a la
vez, porque revisa tanto la forma del documento como su correspondencia con el
dominio del problema.

Entre las técnicas de validación disponibles se distinguen la revisión informal,
el \textit{walkthrough} guiado por el autor, el prototipado para validar
requisitos de interfaz, las listas de verificación basadas en atributos de
calidad, y la revisión por perspectivas, en la que distintos roles (usuario,
desarrollador, probador) inspeccionan el mismo documento desde ángulos diferentes.
De todas ellas, la inspección Fagan es la más formal y la única que exige roles
diferenciados, preparación individual obligatoria y métricas cuantificadas de
defectos \cite{fagan1976}.

El método Fagan, propuesto por Michael Fagan en IBM en 1976, se organiza en fases
secuenciales: planificación (selección del material y de los roles), reunión de
introducción u \textit{overview}, preparación individual, reunión de inspección,
corrección (\textit{rework}) y seguimiento (\textit{follow-up}) \cite{fagan1976}.
Los cuatro roles centrales son el Moderador, que planifica, conduce la reunión y
verifica el cierre de los defectos; el Lector, que parafrasea el documento sección
por sección sin leerlo literalmente, para forzar una segunda interpretación del
texto; y dos Inspectores, que preparan de forma independiente su propia lista de
defectos antes de la reunión. La regla que prohíbe discutir soluciones durante la
reunión de inspección no es arbitraria: busca mantener el foco exclusivo en la
detección, ya que mezclar detección y corrección reduce estadísticamente el número
de defectos encontrados por unidad de tiempo.

Las métricas de inspección más utilizadas son la densidad de defectos (defectos
por página inspeccionada), la distribución por tipo y por severidad, la tasa de
detección por inspector y el esfuerzo total en horas-persona. Estas métricas no
solo cuantifican el resultado de una sesión, sino que permiten comparar la eficacia
de distintos equipos o secciones del documento y decidir si un artefacto requiere
una segunda ronda de inspección antes de avanzar a la siguiente fase del proyecto.
Un caso de estudio recurrente en la literatura es el de los primeros informes
internos de IBM, donde Fagan documentó reducciones sustanciales en la tasa de
defectos que llegaban a las fases de prueba tras la introducción de inspecciones
formales \cite{fagan1976}, resultado que motivó su adopción posterior en procesos
de ingeniería de requisitos y no solo de código.

\subseccion{3.3 Gestión de requisitos}

La gestión de requisitos comprende el control de la configuración del ERS a lo
largo de su evolución: cada versión publicada debe quedar identificada de forma
única, con un historial de revisiones que registre qué cambió, cuándo y por qué
\cite{ieee15288}. Una línea base (\textit{baseline}) es un punto de referencia
formalmente aprobado que congela el estado del documento en un momento dado; a
partir de ella, cualquier modificación posterior debe tramitarse como un cambio
controlado y no como una edición libre. El versionado semántico del ERS (p. ej.,
de v1.0 a v1.1) y el uso de un \textit{tag} anotado en el repositorio Git son la
materialización técnica de ese concepto de línea base.

Cada requisito, además del enunciado, debería llevar atributos de gestión:
identificador único, prioridad, origen o fuente, estado, autor y versión en la
que fue introducido o modificado. Estos atributos son los que permiten construir
la trazabilidad, entendida en dos sentidos: la trazabilidad hacia atrás
(\textit{pre-traceability} o \textit{backward traceability}), que enlaza el
requisito con el stakeholder o documento que le dio origen, y la trazabilidad
hacia adelante (\textit{post-traceability} o \textit{forward traceability}), que
lo enlaza con los artefactos de diseño, código y prueba que lo implementan y
verifican \cite{gotel1994}. La matriz de trazabilidad es la representación tabular
de ambos sentidos y es también la herramienta que sustenta un análisis de impacto
riguroso: ante una solicitud de cambio, el impacto no se estima a criterio del
equipo, sino que se deriva recorriendo la matriz para identificar todos los
artefactos directa e indirectamente afectados \cite{clelandhuang2012}.

El proceso de cambio formal consta de cinco etapas: la solicitud (RFC), que
describe el cambio y su justificación de negocio; el análisis de impacto,
sustentado en la matriz de trazabilidad; la decisión, tomada por el Change
Control Board conforme a lo establecido en el PMBOK \cite{pmbok2021}; la
implementación del cambio aprobado; y el cierre, que actualiza la versión del
documento y su historial. Un indicador complementario es la volatilidad de
requisitos, que mide la proporción de requisitos añadidos, modificados o
eliminados respecto del total en un período dado, y que sirve como señal temprana
de inestabilidad en el alcance del proyecto \cite{wiegers2013}.

\subseccion{3.4 Herramientas CASE para IR}

Las herramientas CASE (\textit{Computer-Aided Software Engineering}) aplicadas a
la ingeniería de requisitos se clasifican habitualmente en tres categorías:
\textit{upper-CASE}, orientadas a las fases tempranas de análisis y
especificación; \textit{lower-CASE}, centradas en construcción, prueba y
despliegue; e \textit{integrated-CASE} (I-CASE), que cubren el ciclo completo y
conectan ambos extremos mediante un repositorio compartido \cite{wiegers2013}.
Jira, Trello, IBM DOORS, Jama Connect y ReqView, comparadas en la sección 9 de
este informe, se ubican principalmente en el segmento upper-CASE o I-CASE según
su grado de integración con herramientas de control de versiones.

Las funcionalidades mínimas esperables de una herramienta CASE de gestión de
requisitos son: un repositorio centralizado con control de acceso; atributos
personalizables por requisito; trazabilidad bidireccional mediante enlaces
explícitos; capacidad de establecer líneas base (\textit{baselining}) sobre
conjuntos de requisitos; un flujo de control de cambios con estados y
aprobaciones; funciones de colaboración como comentarios y notificaciones; e
integración con sistemas de control de versiones (VCS) como Git
\cite{clelandhuang2012}. La selección de una herramienta concreta debe considerar
criterios como el tamaño del equipo, la complejidad regulatoria del dominio, el
presupuesto disponible y la curva de aprendizaje, dado que herramientas más
potentes como DOORS suelen exigir mayor inversión de tiempo de capacitación a
cambio de mayor capacidad de trazabilidad en proyectos regulados.

\subseccion{3.5 IR en procesos ágiles}

En los procesos ágiles, el equivalente funcional del ERS documental es el
\textit{product backlog}: una lista viva y priorizada de historias de usuario,
cada una acompañada de sus criterios de aceptación \cite{cohn2004}. A diferencia
de un requisito formal, una historia de usuario es deliberadamente ligera y actúa
como promesa de una conversación futura entre el equipo y el stakeholder, en la
que se terminan de precisar los detalles. El refinamiento continuo del backlog
(\textit{grooming} o \textit{backlog refinement}) y la definición explícita de
\textit{Definition of Ready} (DoR) y \textit{Definition of Done} (DoD) cumplen,
en un contexto ágil, un papel análogo al de los criterios de verificabilidad y
completitud que la inspección Fagan aplica sobre el ERS documental.

El desarrollo guiado por comportamiento (BDD) formaliza los criterios de
aceptación mediante el formato Gherkin (\textit{Given-When-Then}), que expresa el
comportamiento esperado del sistema en lenguaje semiestructurado, comprensible
tanto por el equipo técnico como por el stakeholder de negocio \cite{smart2014}.
Este formato facilita, además, una forma ligera de trazabilidad: cada escenario
Gherkin puede enlazarse automáticamente con la historia de usuario que lo origina
y con la prueba automatizada que lo verifica, sustituyendo a la matriz de
trazabilidad estática por una trazabilidad ejecutable.

La diferencia principal entre la IR documental y la IR ágil no es de rigor sino
de momento y de granularidad: la primera privilegia una especificación completa y
aprobada antes de construir, mientras que la segunda privilegia la entrega
incremental y difiere el detalle hasta el momento de implementación
\cite{cohn2004}. Ninguno de los dos enfoques es superior de forma absoluta: la IR
documental resulta más adecuada en dominios regulados o con contratos de alcance
fijo, mientras que la IR ágil resulta más eficiente cuando el problema es
exploratorio y el costo de cambio de rumbo debe mantenerse bajo. Esta discusión se
retoma de forma contextualizada, aplicada al SGCV-IA, en la sección 10 de este
informe.

% ================================================================
% 4. INSPECCIÓN FAGAN
% ================================================================
\seccion{4.}{INSPECCIÓN FAGAN}

\subseccion{4.1 Roles y alcance}
\begin{tabular}{ll}
\textbf{Moderador:} & Marcillo Ponce, Alberto Jeanpool \\
\textbf{Lector:} & Arboleda Yanza, Francisco Javier \\
\textbf{Inspector 1:} & Herrera Ramos, Thais Melanie \\
\textbf{Inspector 2:} & Marcillo Ponce, Alberto Jeanpool \\
\end{tabular}

\vspace{0.2cm}
\small\textit{Nota: el equipo de práctica está integrado por tres personas, por lo
que Marcillo Ponce Alberto Jeanpool desempeña simultáneamente los roles de
Moderador e Inspector 2, conforme a lo registrado en su copia del Anexo A.}
\normalsize

\vspace{0.3cm}
\textbf{Alcance inspeccionado:} ERS del SGCV-IA, secciones 3.2 -- Requisitos
Funcionales (RF-01 a RF-25) y 3.3 -- Requisitos No Funcionales (RNF-01 a
RNF-15), equivalentes a 13 páginas efectivas del documento.

\subseccion{4.2 Preparación individual}
Cada inspector completó de forma independiente su propia copia del Anexo A
antes de la reunión de inspección, cubriendo las seis propiedades de calidad de
la ISO/IEC/IEEE 29148:2018 (ambigüedad, completitud, consistencia,
verificabilidad, trazabilidad y factibilidad). Los tres inspectores superaron
el mínimo de 6 defectos exigido por el Paso 1.2 de la guía:

\begin{itemize}
  \item \textbf{Herrera Ramos, Thais Melanie (Inspector 1):} 6 hallazgos,
  concentrados en trazabilidad hacia atrás (RNF-11 a RNF-15 sin RC de origen),
  factibilidad (dependencia de WhatsApp Business API en RF-10/RF-12) y
  consistencia entre RF y su RNF asociado (RF-01/RNF-08).
  \item \textbf{Arboleda Yanza, Francisco Javier (Lector):} 6 hallazgos,
  centrados en consistencia entre requisitos con lógica similar (RF-14 frente
  a RF-17/RF-19; RF-22 frente a RF-24) y en trazabilidad cruzada RF--RNF
  (RF-08/RNF-02).
  \item \textbf{Marcillo Ponce, Alberto Jeanpool (Moderador / Inspector 2):}
  6 hallazgos, centrados en verificabilidad (criterios sin umbral ni
  responsable definido en RF-16, RF-06 y RNF-12) y completitud de casos de
  excepción (RF-21, RF-15, RF-13/RF-14).
\end{itemize}
Las tres copias del Anexo A, firmadas y fechadas de forma individual, se
adjuntan íntegras en la sección de Anexos de este informe.

\subseccion{4.3 Acta de la reunión de inspección}
La reunión de inspección se realizó el 05/08/2026, en modalidad virtual, con
una duración que no superó los 90 minutos establecidos por el Paso 1.3 de la
guía. Asistieron los tres integrantes del equipo en sus roles asignados. El
Lector (Arboleda Yanza, Francisco Javier) parafraseó el ERS sección por
sección --primero el bloque de Requisitos Funcionales (RF-01 a RF-25) y luego
el de Requisitos No Funcionales (RNF-01 a RNF-15)-- sin leerlo literalmente,
conforme al método Fagan. Los dos Inspectores reportaron sus hallazgos en el
orden del recorrido y el Moderador los consolidó en tiempo real en el Anexo B,
asignando tipo, severidad y requisito afectado a cada defecto. Conforme a la
regla del método, no se discutieron soluciones durante la reunión: la
inspección se limitó a la detección; la corrección se documenta por separado
en la sección 5 de este informe.

\subseccion{4.4 Registro consolidado de defectos}
% Mínimo 15 defectos únicos, de los cuales al menos 3 críticos o mayores.
\begin{longtable}{p{1.1cm}p{1.6cm}p{4.6cm}p{1.7cm}p{1.7cm}p{2.6cm}}
\toprule
\textbf{N.\textordmasculine} & \textbf{Sección} & \textbf{Descripción del defecto} & \textbf{Tipo} & \textbf{Severidad} & \textbf{Requisito afectado} \\
\midrule
\endhead
D-01 & 3.3 RNF & RNF-11 a RNF-15 figuran como ``Elaboración propia'', sin RC asociado; rompe la trazabilidad hacia atrás hacia un stakeholder u origen documental. & Trazabilidad & Mayor & RNF-11 -- RNF-15 \\
D-02 & 3.2 RF & RF-10 y RF-12 dependen de WhatsApp Business API sin que el ERS documente el riesgo de dependencia de ese proveedor externo. & Factibilidad & Mayor & RF-10, RF-12 \\
D-03 & 3.2 / 3.3 & RF-01 exige ``la misma funcionalidad esencial'' (100\,\%) entre plataformas; RNF-08 la cuantifica en $\geq$95\,\%, contradiciendo al RF que respalda. & Consistencia & Mayor & RF-01, RNF-08 \\
D-04 & 3.2 RF & El criterio de aceptación de RF-25 menciona un selector de período (día/semana/mes) que no aparece en la descripción del requisito. & Completitud / Consistencia & Menor & RF-25 \\
D-05 & 3.2 RF & RF-17 y RF-19 exigen aprobación explícita del veterinario para sugerencias de IA de diagnóstico; RF-14 (recomendaciones nutricionales, también automáticas) no exige el mismo control. & Consistencia & Mayor & RF-14, RF-17, RF-19 \\
D-06 & 3.2 RF & ``Conforme se aproxime dicha fecha'' no cuantifica los umbrales en el propio texto de RF-05; solo remite implícitamente a RNF-03. & Ambigüedad & Menor & RF-05 \\
D-07 & 3.2 RF & RF-04 no aclara si los diez campos listados en el historial clínico son todos obligatorios o solo un subconjunto. & Ambigüedad & Menor & RF-04 \\
D-08 & 3.2 RF & El criterio ``< 5 segundos'' de RF-16 no define un tamaño de historial de referencia, a diferencia de RNF-01 que sí acota ``hasta 500 pacientes''. & Verificabilidad & Menor & RF-16 \\
D-09 & 3.2 / 3.3 & RF-17 no cita explícitamente a RNF-09 (que agrega el criterio de tiempo $\leq$5 s), rompiendo el patrón de cruce que sí siguen RF-05, RF-08, RF-22 y RF-23. & Trazabilidad & Menor & RF-17, RNF-09 \\
D-10 & 3.2 RF & RF-21 no especifica qué ocurre con los módulos de IA (RF-17--19) cuando el sistema opera sin conexión a internet. & Completitud & Mayor & RF-21 \\
D-11 & 3.2 RF & RF-09 admite ``llamada, WhatsApp, mensaje'' como canales; RF-10 restringe el envío de exámenes solo a WhatsApp sin justificar la diferencia. & Consistencia & Menor & RF-09, RF-10 \\
D-12 & 3.3 RNF & RNF-11 fija 20 usuarios concurrentes sin que ningún RF dimensione el personal esperado de la clínica. & Factibilidad & Menor & RNF-11 \\
D-13 & 3.2 RF & RF-06 no indica quién configura el ``umbral mínimo configurable'' ni el rango permitido. & Verificabilidad & Menor & RF-06 \\
D-14 & 3.2 RF & RF-15 no indica qué ocurre si el recordatorio de indicaciones médicas no se atiende (¿reprogramación?, ¿escalamiento?). & Completitud & Menor & RF-15 \\
D-15 & 3.2 RF & RF-22 y RF-24 regulan ambos el control de acceso por rol; RF-24 remite a RF-22 pero no queda claro si son un mismo requisito fragmentado. & Consistencia & Menor & RF-22, RF-24 \\
D-16 & 3.3 RNF & RNF-12 no define quién audita el cumplimiento de cifrado TLS 1.2/AES-256 ni con qué periodicidad. & Verificabilidad & Menor & RNF-12 \\
D-17 & 3.2 / 3.3 & RF-08 no cita a RNF-02 en su propia descripción (solo aparece en el criterio de aceptación), a diferencia de otros pares RF/RNF que sí se citan mutuamente. & Trazabilidad & Menor & RF-08, RNF-02 \\
D-18 & 3.2 RF & RF-14 depende del dato de peso/condición corporal (RF-13) pero no define qué ocurre si aún no existe historial suficiente (primera visita). & Completitud & Mayor & RF-13, RF-14 \\
\bottomrule
\end{longtable}

\textit{Detectado por:} D-04, D-05, D-07, D-11, D-15, D-17 -- Arboleda Yanza,
Francisco Javier (Lector); D-01, D-02, D-03, D-06, D-09, D-12 -- Herrera Ramos,
Thais Melanie (Inspector 1); D-08, D-10, D-13, D-14, D-16, D-18 -- Marcillo
Ponce, Alberto Jeanpool (Moderador / Inspector 2).

\subseccion{4.5 Métricas de inspección}

\begin{longtable}{p{3.6cm}p{2.4cm}p{7.2cm}}
\toprule
\textbf{Métrica} & \textbf{Valor} & \textbf{Interpretación} \\
\midrule
\endhead
Densidad de defectos (def./pág.) & 1{,}38 (18/13) & Con 18 defectos sobre 13 páginas
efectivas, el ERS presenta una densidad moderada: los requisitos están
mayormente bien redactados de forma aislada, pero persisten fallas que solo
emergen al comparar requisitos entre sí. \\
\addlinespace
Distribución por tipo & Consistencia (5), Completitud (4), Trazabilidad (3),
Verificabilidad (3), Ambigüedad (2), Factibilidad (2) & Los defectos de
consistencia y completitud concentran la mayoría de los hallazgos, lo que
indica que el ERS necesita revisión cruzada entre requisitos relacionados
(especialmente entre un RF y su RNF asociado) más que corrección de redacción
aislada. \\
\addlinespace
Distribución por severidad & Crítico: 0; Mayor: 6 (33\,\%); Menor: 12 (67\,\%) &
Ningún defecto detectado es crítico, pero una tercera parte son mayores y
exigen corrección obligatoria antes de establecer la línea base, conforme al
Paso 1.5 de la guía. \\
\addlinespace
Tasa de detección por inspector & 6 defectos cada uno (Thais, Francisco,
Jeanpool) & La distribución equilibrada (6/6/6) evidencia una preparación
individual proporcional entre los tres integrantes, sin concentración de
hallazgos en un solo inspector. \\
\addlinespace
Esfuerzo total (horas-persona) & $\approx$ 3,5 h-p (3 integrantes $\times$
$\sim$70 min de preparación individual y reunión de 90 min) & El esfuerzo se
mantuvo dentro de las 2 horas asignadas al Paso 1 de la guía para un ERS de 13
páginas efectivas. \\
\bottomrule
\end{longtable}

% ================================================================
% 5. CORRECCIONES APLICADAS
% ================================================================
\seccion{5.}{CORRECCIONES APLICADAS}

% Tabla antes/después por cada defecto crítico y mayor. 100% deben corregirse.
\begin{longtable}{p{0.9cm}p{4.3cm}p{4.6cm}p{2.1cm}}
\toprule
\textbf{ID} & \textbf{Texto antes} & \textbf{Texto después} & \textbf{Requisito modificado} \\
\midrule
\endhead
D-01 &
Columna ``RC relacionado'' = ``Elaboración propia'' (sin vínculo a un RC de origen) en RNF-11 a RNF-15. &
Cada RNF queda enlazado a un stakeholder real del Mapa de Stakeholders: RNF-11 $\to$ Equipo de desarrollo; RNF-12 $\to$ Comité de ética/protección de datos (LOPDP); RNF-13 $\to$ Propietario/administrador de la clínica; RNF-14 $\to$ Equipo de desarrollo; RNF-15 $\to$ Personal administrativo y Médico veterinario. &
RNF-11 -- RNF-15 \\
\addlinespace
D-02 &
``El sistema debe permitir adjuntar y enviar resultados de exámenes directamente por WhatsApp desde la plataforma.'' (RF-10) &
Se agrega: ``Ante la indisponibilidad del proveedor de mensajería (WhatsApp Business API), el sistema debe reintentar el envío y notificar al personal administrativo del fallo.'' Este comportamiento se formaliza como nuevo requisito RF-26 vía RFC-02 (sección 6.2). &
RF-10, RF-12 (nuevo RF-26) \\
\addlinespace
D-03 &
RF-01: ``...manteniendo la misma funcionalidad esencial en ambos entornos.'' / RNF-08: ``...con paridad de funciones $\geq$95\,\% entre plataformas.'' &
RF-01 (corregido): ``...manteniendo una paridad de funcionalidad esencial de al menos 95\,\% entre ambos entornos, conforme a RNF-08.'' &
RF-01 \\
\addlinespace
D-05 &
``El sistema debe generar recomendaciones de alimentación a partir de la raza, edad, peso y condición de salud del animal, como apoyo al juicio del veterinario.'' &
``...sin que ninguna recomendación pueda aplicarse sin revisión explícita del veterinario, como apoyo a su juicio clínico.'' &
RF-14 \\
\addlinespace
D-10 &
``El sistema debe permitir el registro y la consulta de historiales clínicos de manera local sin conexión a internet, sincronizando al restablecerse la conexión.'' &
Se agrega: ``Los módulos dependientes de IA (RF-17--RF-19) quedan deshabilitados durante la operación sin conexión y se reactivan automáticamente al restablecerse la conectividad, sin afectar el registro y consulta de historiales clínicos.'' &
RF-21 \\
\addlinespace
D-18 &
``La recomendación se actualiza automáticamente al registrarse un nuevo dato de peso o condición corporal (RF-13).'' &
Se agrega: ``En la primera visita, cuando aún no exista historial suficiente, el sistema debe notificar la ausencia de datos y omitir la generación de la recomendación hasta contar con al menos un registro histórico.'' &
RF-14, RF-13 \\
\bottomrule
\end{longtable}

\vspace{0.2cm}
\textbf{Defectos menores no corregidos y su justificación:} conforme al Paso
1.5 de la guía, los 12 defectos menores (D-04, D-06, D-07, D-08, D-09, D-11,
D-12, D-13, D-14, D-15, D-16, D-17) no comprometen la validez funcional del
ERS ni bloquean la construcción de la línea base v1.1; quedan documentados
como mejoras de redacción y de trazabilidad fina pendientes para la siguiente
revisión del documento (Entrega 2B), en lugar de tramitarse como RFC en esta
práctica. Dos de ellos (D-04 y D-18) están parcialmente cubiertos por las
correcciones de D-05 y D-10 respectivamente, al compartir el mismo requisito
afectado.

% ================================================================
% 6. GESTIÓN DEL CAMBIO Y CCB
% ================================================================
\seccion{6.}{GESTIÓN DEL CAMBIO Y CCB}

\subseccion{6.1 Solicitud de Cambio RFC-01 -- (Alcance / nuevo requisito)}
\begin{longtable}{p{4.3cm}p{8.5cm}}
\toprule
\multicolumn{2}{l}{\tituloTabla{RFC-01 -- 05/08/2026 -- Calidad (modificación de NFR)}} \\
\midrule
\endhead
ID / Fecha / Solicitante & RFC-01 -- 05/08/2026 -- Herrera Ramos, Thais Melanie (Inspector 1) \\
Requisito afectado & RF-01, RNF-08 \\
Descripción del cambio & Detectado en la inspección Fagan (defecto D-03, Anexo B): RF-01 exige ``la misma funcionalidad esencial'' (100\,\%) entre plataformas móvil y de escritorio, mientras que RNF-08 cuantifica la paridad en $\geq$95\,\%. Se solicita modificar RF-01 para que remita explícitamente a la métrica de RNF-08, eliminando la contradicción. \\
Justificación de negocio & Un RF y su RNF asociado no pueden exigir umbrales distintos para el mismo atributo de calidad; mantener ambos textos vigentes genera ambigüedad para desarrollo y para la verificación de aceptación. \\
Requisitos impactados (matriz) & RF-01 (texto modificado); RNF-08 (sin cambio, queda como referencia); indirectamente RF-20 (usabilidad multiplataforma), por compartir el criterio de paridad funcional. \\
Artefactos impactados & Documento ERS (secciones 3.2 y 3.3); casos de prueba de aceptación de RF-01; mockups de interfaz móvil/escritorio. \\
Costo estimado (h-p) / Riesgo & 1 hora-persona / Riesgo bajo: corrección de consistencia textual que no altera comportamiento ni esfuerzo de desarrollo ya planificado. \\
Recomendación técnica & Aprobar sin condiciones: alinea el ERS consigo mismo sin alterar el alcance funcional acordado con los stakeholders. \\
\bottomrule
\end{longtable}

\subseccion{6.2 Solicitud de Cambio RFC-02 -- (Alcance / nuevo requisito)}
\begin{longtable}{p{4.3cm}p{8.5cm}}
\toprule
\multicolumn{2}{l}{\tituloTabla{RFC-02 -- 05/08/2026 -- Alcance (nuevo requisito)}} \\
\midrule
\endhead
ID / Fecha / Solicitante & RFC-02 -- 05/08/2026 -- Arboleda Yanza, Francisco Javier (Lector) \\
Requisito afectado & RF-10, RF-12 (nuevo: RF-26) \\
Descripción del cambio & La inspección Fagan (defecto D-02, Anexo B) identificó que RF-10 y RF-12 dependen de WhatsApp Business API sin mecanismo de contingencia ante indisponibilidad del proveedor. Se solicita incorporar RF-26: el sistema debe reintentar el envío y notificar al personal administrativo cuando el canal de mensajería falle. \\
Justificación de negocio & El Mapa de Stakeholders ya identifica al ``Proveedor del servicio de mensajería'' como stakeholder de bajo poder/bajo interés bajo estrategia de ``monitorear''; su falla no debe paralizar el envío de resultados de exámenes ni de recordatorios de citas. \\
Requisitos impactados (matriz) & RF-10, RF-12 (quedan referenciados desde RF-26); RNF-04 (disponibilidad) y RNF-10 (modularidad), por tratarse también de tolerancia a fallos de un servicio externo. \\
Artefactos impactados & Documento ERS (nueva entrada RF-26 en 3.2); matriz de trazabilidad (nueva fila); backlog en la herramienta CASE (nueva Story). \\
Costo estimado (h-p) / Riesgo & 3 horas-persona / Riesgo medio: sin este RF, el sistema queda sin comportamiento definido ante falla de un proveedor externo. \\
Recomendación técnica & Aprobar con la condición de cuantificar el número de reintentos y el tiempo máximo de espera antes de notificar la falla. \\
\bottomrule
\end{longtable}

\subseccion{6.3 Solicitud de Cambio RFC-03 -- (Restricción / normativa)}
\begin{longtable}{p{4.3cm}p{8.5cm}}
\toprule
\multicolumn{2}{l}{\tituloTabla{RFC-03 -- 05/08/2026 -- Restricción / normativa}} \\
\midrule
\endhead
ID / Fecha / Solicitante & RFC-03 -- 05/08/2026 -- Marcillo Ponce, Alberto Jeanpool (Moderador) \\
Requisito afectado & RNF-12 (nuevo: RNF-16) \\
Descripción del cambio & Se solicita incorporar RNF-16, que establezca el derecho de propietarios de mascotas y personal a solicitar rectificación o eliminación de sus datos personales, conforme a la Ley Orgánica de Protección de Datos Personales (LOPDP) de Ecuador. \\
Justificación de negocio & El Mapa de Stakeholders identifica al ``Comité de ética / protección de datos (LOPDP)'' como stakeholder de alto poder bajo estrategia ``gestionar de cerca''. El ERS actual regula cifrado y auditoría (RNF-12, RNF-06) pero no los derechos ARCO (acceso, rectificación, cancelación, oposición). \\
Requisitos impactados (matriz) & RNF-06 (trazabilidad/no repudio) y RNF-12 (cifrado), con los que RNF-16 debe ser consistente; RF-02 (centralización de historiales), sobre el que recaería técnicamente la eliminación solicitada. \\
Artefactos impactados & Documento ERS (nueva entrada RNF-16 en 3.3); matriz Ley-Artículo $\to$ RF/RNF elaborada por el equipo para el criterio legal. \\
Costo estimado (h-p) / Riesgo & 4 horas-persona / Riesgo alto si se rechaza: un sistema que gestiona historiales clínicos sin mecanismo de rectificación/eliminación queda expuesto a incumplimiento normativo. \\
Recomendación técnica & Aprobar. Es una obligación normativa ya anticipada por el trabajo de trazabilidad legal del equipo. \\
\bottomrule
\end{longtable}

\subseccion{6.4 Acta del CCB}
\textbf{Fecha:} 05/08/2026 \quad \textbf{Modalidad:} Virtual

\textbf{Asistentes (roles):} Marcillo Ponce, Alberto Jeanpool -- Presidente del
CCB; Herrera Ramos, Thais Melanie -- Representante del cliente; Arboleda
Yanza, Francisco Javier -- Analista de requisitos / Desarrollador.

\textbf{Agenda:} 1) Revisión de RFC-01 (calidad -- RF-01/RNF-08). 2) Revisión
de RFC-02 (alcance -- nuevo RF-26). 3) Revisión de RFC-03
(restricción/normativa -- nuevo RNF-16). 4) Cierre y actualización de versión
del ERS.

\begin{longtable}{p{1.7cm}p{5.4cm}p{2.1cm}p{4.0cm}}
\toprule
\textbf{RFC} & \textbf{Deliberación resumida} & \textbf{Decisión} & \textbf{Justificación / Responsable / Plazo} \\
\midrule
\endhead
RFC-01 & El Analista señala que la corrección no altera el alcance funcional acordado; el Representante del cliente confirma que el 95\,\% de paridad es aceptable, ya que las funciones críticas (historial, citas, facturación) están cubiertas al 100\,\% en ambos entornos. Sin objeciones. &
Aprobado (sin condiciones) &
No requiere justificación adicional -- consenso directo. Responsable: Herrera Ramos, Thais Melanie. Plazo: misma sesión de trabajo. \\
\addlinespace
RFC-02 & El Presidente considera razonables las 3 h-p; el Analista advierte que el criterio original no cuantificaba el mecanismo de reintento, dejando el requisito tan ambiguo como el defecto que corrige. El Representante del cliente propone un límite concreto. &
Aprobado con condiciones &
Debe cuantificarse: máximo 3 reintentos automáticos, intervalo de 5 min entre cada uno, notificación al personal administrativo si no se logra el envío en 20 min. Responsable: Arboleda Yanza, Francisco Javier. Plazo: antes de la línea base v1.1. \\
\addlinespace
RFC-03 & El Representante del cliente señala que eliminar historiales clínicos por completo podría chocar con obligaciones de conservación sanitaria. El Analista propone distinguir datos administrativos/de contacto (eliminables) de historiales clínicos (con conservación mínima obligatoria). El Presidente acoge la propuesta. &
Aprobado con condiciones &
Debe distinguirse explícitamente en RNF-16 entre datos administrativos e historiales clínicos, para no crear conflicto con RNF-06. Responsable: Marcillo Ponce, Alberto Jeanpool. Plazo: antes de la línea base v1.1. \\
\bottomrule
\end{longtable}
% Decisión ∈ {aprobado, aprobado con condiciones, rechazado, diferido}

\textbf{Versión resultante del ERS:} de v1.0 (Entrega 2A) a v1.1, que
incorpora la corrección de RF-01, el nuevo requisito RF-26 y el nuevo
requisito RNF-16. Esta versión debe declararse de forma idéntica en la
carátula del informe, en el historial de revisiones del ERS, en
\texttt{CHANGELOG.md} y en el tag anotado de Git (\texttt{baseline-v1.1}),
conforme al criterio de piso G2 y al criterio C5 de la rúbrica. \\
\textbf{Firmas:} Marcillo Ponce, Alberto Jeanpool; Herrera Ramos, Thais
Melanie; Arboleda Yanza, Francisco Javier \emph{(firmas autógrafas u
ológrafas pendientes de digitalizar en el Acta original, Anexo C).}

% ================================================================
% 7. TRAZABILIDAD EN HERRAMIENTA CASE
% ================================================================
\seccion{7.}{TRAZABILIDAD EN HERRAMIENTA CASE}

\subseccion{7.1 Estructura del tablero}
\textbf{Herramienta utilizada:} Jira (proyecto ``SI'' / ``Y'', cuenta gratuita
del equipo). \\[4pt]
\textcolor{uteqgris}{\textbf{Pendiente de cierre por el equipo antes de compilar
la versión final:} el tablero creado hasta el momento organiza el backlog con
una numeración simplificada (RF-01 a RF-15, un issue por Epic funcional:
Gestión clínica, Inteligencia artificial y seguimiento, Gestión
administrativa, Inventario y seguridad) que \textbf{no coincide literalmente}
con la numeración oficial del ERS vigente (RF-01 a RF-25, según
\texttt{SGCVIA\_RF\_RNF\_Formato\_ERS.docx}). Por ejemplo, el tablero actual
llama ``RF-01'' a ``Registrar historial clínico'', mientras que en el ERS
RF-01 es ``Acceso multiplataforma'' y el registro del historial corresponde a
RF-04. Antes de exportar el CSV definitivo (Anexo D) y de construir la matriz
de 7.3, el equipo debe renombrar las Stories del tablero para que su ID y
título coincidan exactamente con el RF del ERS que representan; de lo
contrario, la matriz de trazabilidad quedaría construida sobre requisitos que
no existen tal como están nombrados en el ERS, lo que activa el criterio de
severidad S3 de la guía (matriz con elementos inexistentes en el proyecto).
Ver la comunicación adjunta al equipo con el detalle de esta corrección.}

\subseccion{7.2 Campos personalizados}
El equipo intentó configurar los campos ``Prioridad MoSCoW'' y ``Fuente del
requisito'' en Jira, sin éxito por restricciones de permisos del plan
gratuito del sitio (no se encontró la opción de creación de campos
personalizados a nivel de proyecto). \textcolor{uteqgris}{\textbf{Pendiente:}
si el sitio de Jira permite finalmente crear estos campos con un
administrador del sitio, deben configurarse y aplicarse a las 15+ Stories; si
no es posible, el equipo debe documentar aquí esta limitación real (tal como
se hizo) y suplir la información de MoSCoW y Fuente del requisito mediante una
columna adicional en la matriz de 7.3 y en el CSV exportado, dejando
constancia de la limitación de la herramienta -- esto es preferible a dejar
la sección vacía, conforme al criterio C4 de la rúbrica (``tablero parcial''
se califica en desarrollo, no ausente, si la limitación está documentada).}

\subseccion{7.3 Matriz de trazabilidad}
% Mínimo 15 filas: Stakeholder -> RF -> CU -> Clase/Módulo -> Prueba
\textcolor{uteqgris}{\textbf{Pendiente:} esta matriz debe construirse
\emph{después} de corregir la numeración del tablero (7.1). Sus columnas
deben usar el RF real del ERS (RF-01 a RF-25) y el stakeholder debe tomarse
del Mapa de Stakeholders del equipo (\texttt{Mapa\_de\_stakeholders.csv}), ya
usado de forma consistente en las secciones 5 y 6 de este informe. La columna
CU/Clase/Módulo solo puede completarse con nombres reales: si el equipo no
cuenta todavía con un diagrama de casos de uso o de clases del PFC, debe
consignarse el módulo del sistema (p. ej. ``Módulo Clínico'', ``Módulo IA'',
``Módulo Administrativo'', ``Módulo Inventario/Seguridad'') en lugar de
inventar un CU o una clase, y dejar esa limitación explícita en esta misma
subsección, tal como permite el criterio C4.}
\begin{longtable}{p{2.3cm}p{1.6cm}p{1.6cm}p{2.8cm}p{2.5cm}}
\toprule
\textbf{Stakeholder} & \textbf{RF} & \textbf{CU} & \textbf{Clase / Módulo} & \textbf{Prueba} \\
\midrule
\endhead
& & & & \\
% [PENDIENTE: mínimo 15 filas, una por cada RF real del ERS, con su stakeholder de origen]
\bottomrule
\end{longtable}

\textbf{Requisitos huérfanos identificados:} [PENDIENTE: listar los RF/RNF del
ERS que aún no tengan Story ni enlace de trazabilidad en el tablero, una vez
corregida la numeración.]

\subseccion{7.4 Evidencias}
[PENDIENTE: referenciar capturas del tablero completo, detalle de \(\geq 3\)
issues con enlaces visibles y panel de estadísticas -- ver Anexo E. Referenciar
exportación CSV -- ver Anexo D. Estas capturas deben tomarse \emph{después} de
la corrección de numeración de 7.1 para ser coherentes con el resto del
informe.]

% ================================================================
% 8. LÍNEA BASE CON GIT
% ================================================================
\seccion{8.}{LÍNEA BASE CON GIT}

\subseccion{8.1 Repositorio}
\textbf{URL:} \url{https://github.com/[NombreEquipo]/ISR401-PFC-ERS-SGCVIA}
\textcolor{uteqgris}{\textit{[PENDIENTE: reemplazar por la URL pública real
una vez creado el repositorio -- criterio de piso G1: si esta URL no es
accesible públicamente al momento de la entrega, la calificación es CERO.]}}

\subseccion{8.2 Commits semánticos}
[PENDIENTE: el equipo debe registrar al menos 8 commits distribuidos, con
autoría de los tres integrantes, documentando la subida del ERS v1.0, la
aplicación de cada uno de los tres cambios aprobados (RFC-01, RFC-02, RFC-03)
y la publicación de v1.1. Ejemplo de mensaje semántico coherente con esta
práctica: \texttt{docs(ers): v1.1 - aplicar RFC-01/02/03 aprobados por CCB
semana 14}.]

\subseccion{8.3 Tag de línea base}
\begin{verbatim}
git tag -a baseline-v1.1 -m "Baseline aprobada por CCB - RFC-01/02/03"
git push origin baseline-v1.1
\end{verbatim}
Esta versión (v1.1) debe coincidir exactamente con la declarada en el Acta
del CCB (sección 6.4) y en el \texttt{CHANGELOG.md} (8.4).

\subseccion{8.4 CHANGELOG.md}
[PENDIENTE: una vez aplicados los cambios en el repositorio, transcribir aquí
el \texttt{CHANGELOG.md} real. Estructura sugerida, ya coherente con las
decisiones del CCB de la sección 6.4:]
\begin{verbatim}
## [1.1] - 2026-08-05
### Añadido
- RF-26: reintento y notificación ante falla del proveedor de mensajería
  (RFC-02, aprobado con condiciones).
- RNF-16: derechos de rectificación/eliminación de datos personales
  conforme a LOPDP (RFC-03, aprobado con condiciones).
### Cambiado
- RF-01: paridad de funcionalidad esencial ajustada a >= 95%,
  alineada con RNF-08 (RFC-01, aprobado sin condiciones).
### Eliminado
- (ninguno en esta línea base)
\end{verbatim}

\subseccion{8.5 Historial (ver Anexo E)}
[PENDIENTE: referenciar capturas de \texttt{git log --oneline --graph
--decorate} y \texttt{git tag -n}, tomadas después de publicar el tag
\texttt{baseline-v1.1}.]

% ================================================================
% 9. COMPARATIVA DE HERRAMIENTAS CASE
% ================================================================
\seccion{9.}{COMPARATIVA DE HERRAMIENTAS CASE}

Se comparan cuatro herramientas CASE de gestión de requisitos --Jira, IBM DOORS
Next, Jama Connect y ReqView-- bajo seis criterios. Los juicios de cada celda
son propios del equipo, contrastados con documentación oficial y con el uso real
que el equipo dio a Jira/Trello en el Paso 3 de esta práctica; no son una copia
del sitio comercial de cada proveedor.

{\footnotesize
\setlength{\tabcolsep}{3pt}
\begin{longtable}{p{1.7cm}p{2.1cm}p{2.1cm}p{2.1cm}p{2.1cm}p{1.7cm}p{1.9cm}}
\toprule
\textbf{Herramienta} & \textbf{Func. de IR} & \textbf{Trazabilidad} & \textbf{Colaboración} & \textbf{Integr. VCS} & \textbf{Costo} & \textbf{Curva de aprendizaje} \\
\midrule
\endhead
Jira &
Backlog, épicas e historias; sin matriz de trazabilidad nativa, requiere campos personalizados o plugins &
Manual, mediante enlaces y campos; suficiente para un equipo académico pequeño &
Alta; tablero compartido, comentarios y notificaciones en tiempo real &
Nativa con Bitbucket/GitHub vía plugins gratuitos &
Gratuito hasta 10 usuarios; plan pago para equipos mayores &
Baja; interfaz intuitiva, apta sin capacitación formal \\
\addlinespace
IBM DOORS Next &
Especializada en RM formal; atributos y líneas base nativas &
Trazabilidad viva end-to-end, la más robusta de las cuatro para entornos regulados &
Media; pensada para equipos grandes dentro del ecosistema IBM Jazz &
Integración vía conectores con Jira y Git, no nativa &
Alto; licenciamiento empresarial fuera de alcance de un curso &
Alta; exige capacitación específica antes de ser productivo \\
\addlinespace
Jama Connect &
RM moderna centrada en \textit{Live Traceability}; fuerte en análisis de impacto &
Muy alta; visibilidad en tiempo real de relaciones requisito-prueba-riesgo &
Alta; diseñada para colaboración entre stakeholders técnicos y de negocio &
Integra con Jira y Git mediante conectores oficiales &
Alto; orientado a industrias reguladas (automotriz, médico, aeroespacial) &
Media; interfaz más accesible que DOORS pero con curva propia \\
\addlinespace
ReqView &
Organiza requisitos, riesgos y pruebas en un mismo espacio; ligera y enfocada &
Alta; trazabilidad end-to-end declarada como su función distintiva &
Baja--media; pensada para equipos pequeños, menos colaborativa que Jama &
Soporta ReqIF e integración con Git y Jira &
Medio; más accesible que DOORS o Jama para equipos pequeños &
Media; requiere aprender el modelo de datos propio \\
\bottomrule
\end{longtable}
}

\textbf{Recomendación para el PFC del equipo:} para el SGCV-IA --un sistema
académico, sin exigencias regulatorias de la magnitud de un dominio médico o
aeroespacial, desarrollado por un equipo de 3 a 4 integrantes-- Jira es la
opción más adecuada. Su costo nulo en el nivel usado por el equipo, su
integración directa con GitHub y su baja curva de aprendizaje pesan más que la
trazabilidad nativa superior de Jama Connect o DOORS Next, cuyo costo de
licenciamiento y complejidad de configuración no se justifican para un proyecto
académico de esta escala. Si el SGCV-IA evolucionara hacia un despliegue
real con componentes de decisión clínica asistida por IA sujetos a mayor
escrutinio, Jama Connect sería la alternativa a evaluar en primer lugar, por su
trazabilidad viva y su orientación a dominios regulados.

% ================================================================
% 10. IR TRADICIONAL FRENTE A IR ÁGIL
% ================================================================
\seccion{10.}{IR TRADICIONAL FRENTE A IR ÁGIL}

\begin{longtable}{p{2.6cm}p{4.7cm}p{4.7cm}}
\toprule
\textbf{Dimensión} & \textbf{IR tradicional} & \textbf{IR ágil} \\
\midrule
\endhead
Documentación &
ERS completo, aprobado y baselineado antes de iniciar la construcción; cambios
posteriores tramitados formalmente &
Backlog vivo de historias de usuario; el detalle se completa justo antes de
implementarse (\textit{just-in-time}) \\
\addlinespace
Gestión del cambio &
Proceso formal vía RFC y CCB, con análisis de impacto sobre la matriz de
trazabilidad antes de aprobar &
El cambio se absorbe en la reprioritización del backlog entre sprints; no
requiere un comité formal \\
\addlinespace
Participación de stakeholders &
Concentrada en las fases de elicitación y validación (p. ej., la inspección
Fagan de esta práctica) &
Continua, mediante revisión de sprint y refinamiento periódico del backlog \\
\addlinespace
Herramientas &
Herramientas CASE de RM formal (DOORS, Jama) con matrices y líneas base &
Tableros ágiles (Jira, Trello) centrados en flujo de trabajo más que en
documentación normativa \\
\addlinespace
Velocidad &
Más lenta al inicio por el esfuerzo de especificación completa; más predecible
después &
Entrega incremental más rápida; el costo de los cambios se paga en cada sprint \\
\addlinespace
Escalabilidad &
Escala bien en proyectos grandes, multiequipo o con exigencias de auditoría y
cumplimiento normativo &
Escala con dificultad sin prácticas adicionales (p. ej. SAFe) cuando crecen el
número de equipos y las dependencias \\
\bottomrule
\end{longtable}

\textbf{Discusión contextual:} el SGCV-IA, al ser un proyecto académico de
alcance acotado y equipo reducido, se beneficia hoy de un enfoque
predominantemente documental --que es, de hecho, el exigido por esta unidad y
por el curso-- porque permite practicar formalmente la validación, la
trazabilidad y el control de cambios sobre un ERS completo. Sin embargo, los
mecanismos ágiles (historias de usuario con criterios de aceptación en formato
Gherkin, backlog priorizado) resultarían más convenientes si el equipo pasara a
una fase de implementación iterativa del sistema, donde la velocidad de entrega
y la retroalimentación continua del cliente simulado pesan más que la
formalidad del proceso de cambio. En la práctica, ambos enfoques no son
mutuamente excluyentes: la matriz de trazabilidad construida en la sección 7 de
este informe puede alimentar directamente un backlog ágil si el equipo decide
continuar el desarrollo del SGCV-IA bajo Scrum en fases posteriores del PFC.

% ================================================================
% 11. CONCLUSIONES CRÍTICAS
% ================================================================
\seccion{11.}{CONCLUSIONES CRÍTICAS}

% Cinco conclusiones que integren los cuatro temas de la unidad (inspección, CCB,
% trazabilidad, baseline/comparativa). Deben interpretar resultados propios, no
% ser afirmaciones generales de manual.
\begin{enumerate}
  \item De los 18 defectos consolidados en la inspección, 5 de los 6 mayores
  (D-01, D-03, D-05, D-10, D-18) son defectos de \emph{consistencia y
  completitud cruzada} entre un RF y su RNF asociado o entre dos RF de lógica
  similar, no errores de redacción aislados; esto indica que el mayor riesgo
  de calidad del ERS del SGCV-IA no está en cómo se escribe cada requisito por
  separado, sino en la falta de una revisión sistemática de coherencia entre
  requisitos relacionados, algo que una lectura individual (sin roles
  diferenciados) difícilmente habría detectado.
  \item Dos de los tres RFC aprobados (RFC-02 y RFC-03) surgieron de vacíos
  que el propio Mapa de Stakeholders del equipo ya anticipaba --la
  dependencia del proveedor de mensajería y la exigencia del Comité de
  ética/LOPDP--, lo que muestra que la matriz de partes interesadas no es un
  artefacto decorativo: usarla activamente durante el análisis de impacto
  permitió al CCB fundamentar el riesgo de cada solicitud (medio y alto,
  respectivamente) en un stakeholder identificado y no en una estimación
  ``a ojo''.
  \item La condición que el CCB impuso a RFC-02 (cuantificar reintentos y
  tiempo de espera) y a RFC-03 (distinguir datos administrativos de
  historiales clínicos) confirma que aprobar un cambio sin exigir que su
  criterio de aceptación quede tan cuantificado como el resto del ERS
  simplemente traslada el defecto original --la falta de verificabilidad-- del
  requisito viejo al nuevo; el CCB actuó, en la práctica, como un segundo
  filtro de calidad y no solo como un órgano de aprobación de alcance.
  \item El defecto D-01 (RNF-11 a RNF-15 sin RC de origen) fue, a la vez, el
  hallazgo más frecuente entre los tres inspectores y el más costoso de
  corregir en términos de trazabilidad, porque exigió recorrer el Mapa de
  Stakeholders completo para asignar un origen documental verificable a cada
  RNF; esto evidencia que los requisitos no funcionales de ``elaboración
  propia'' son, en este proyecto, el punto más débil de la trazabilidad hacia
  atrás, y deberían ser la prioridad de revisión en la próxima entrega del ERS.
  \item Al construir la sección 7 de este informe, el equipo detectó que el
  tablero CASE (Jira) fue poblado con una numeración de RF simplificada que no
  corresponde exactamente a la numeración oficial del ERS vigente; este
  hallazgo, surgido de la propia práctica y no de la guía, confirma en la
  práctica lo que la sección 3.4 describe en teoría: una herramienta CASE solo
  aporta trazabilidad real si el equipo mantiene disciplina de sincronización
  entre el artefacto normativo (el ERS) y el artefacto operativo (el
  backlog); de lo contrario, la trazabilidad que muestra el tablero es
  ilusoria y debe corregirse antes de exportarse como evidencia.
\end{enumerate}

% ================================================================
% 12. DISTRIBUCIÓN DEL TRABAJO
% ================================================================
\seccion{12.}{DISTRIBUCIÓN DEL TRABAJO}

% Debe ser contrastable con el historial de commits.
\begin{longtable}{p{3.5cm}p{6.0cm}p{3.0cm}}
\toprule
\textbf{Integrante} & \textbf{Aportes} & \textbf{Rol / Commits} \\
\midrule
\endhead
Marcillo Ponce, Alberto Jeanpool &
Moderador de la inspección Fagan (planificación, conducción y consolidación
del Anexo B); 6 defectos detectados como Inspector 2 (D-08, D-10, D-13, D-14,
D-16, D-18); Presidente del CCB; solicitante y responsable de RFC-03
(RNF-16). &
Moderador / Inspector 2; Presidente CCB. Commits: \textit{[PENDIENTE: completar
con el historial real de GitHub]} \\
\addlinespace
Herrera Ramos, Thais Melanie &
6 defectos detectados como Inspector 1 (D-01, D-02, D-03, D-06, D-09, D-12),
incluido el defecto que originó RFC-01; Representante del cliente en el CCB;
solicitante y responsable de RFC-01 (RF-01/RNF-08). &
Inspector 1; Representante del cliente CCB. Commits: \textit{[PENDIENTE]} \\
\addlinespace
Arboleda Yanza, Francisco Javier &
Lector de la reunión de inspección (recorrido sección por sección del ERS);
6 defectos detectados (D-04, D-05, D-07, D-11, D-15, D-17); Analista de
requisitos / Desarrollador en el CCB; solicitante y responsable de RFC-02
(RF-10/RF-12, nuevo RF-26). &
Lector; Analista CCB. Commits: \textit{[PENDIENTE]} \\
\bottomrule
\end{longtable}

\textcolor{uteqgris}{\textit{[PENDIENTE: completar la columna de commits con el
autor real de cada commit según \texttt{git shortlog -sne} una vez publicado
el repositorio, para que esta tabla sea contrastable con el historial de Git
conforme al criterio C7 de la rúbrica.]}}

% ================================================================
% 13. REFERENCIAS
% ================================================================
\seccion{13.}{REFERENCIAS}

\renewcommand{\refname}{\vspace{-2.2em}}
\begin{thebibliography}{99}

\bibitem{ieee29148} IEEE/ISO/IEC, \textit{29148:2018 -- Systems and Software
Engineering: Life Cycle Processes: Requirements Engineering}. IEEE, 2018.
doi: 10.1109/IEEESTD.2018.8559686.

\bibitem{ieee15288} IEEE/ISO/IEC, \textit{15288:2023 -- Systems and Software
Engineering: System Life Cycle Processes}. IEEE, 2023.

\bibitem{ieee12207} IEEE/ISO/IEC, \textit{12207:2017 -- Systems and Software
Engineering: Software Life Cycle Processes}. IEEE, 2017.

\bibitem{iso25010} ISO/IEC, \textit{25010:2023 -- Systems and Software Quality
Requirements and Evaluation (SQuaRE): Product Quality Model}. ISO, 2023.

\bibitem{pohl2025} K. Pohl, \textit{Requirements Engineering: Fundamentals,
Principles, and Techniques}, 2nd ed. Berlin, Germany: Springer, 2025.
doi: 10.1007/978-3-662-69204-2.

\bibitem{swebok2024} H. Washizaki (Ed.), \textit{Guide to the Software
Engineering Body of Knowledge (SWEBOK Guide)}, Version 4.0. Los Alamitos, CA,
USA: IEEE Computer Society, 2024.

\bibitem{pmbok2021} Project Management Institute, \textit{A Guide to the
Project Management Body of Knowledge (PMBOK Guide)}, 7th ed. Newtown Square,
PA, USA: PMI, 2021.

\bibitem{ireb2022} International Requirements Engineering Board, \textit{IREB
Certified Professional for Requirements Engineering -- Foundation Level
Syllabus}, Version 3.1. Karlsruhe, Germany: IREB, 2022.

\bibitem{fagan1976} M. E. Fagan, ``Design and code inspections to reduce
errors in program development,'' \textit{IBM Systems Journal}, vol. 15,
no. 3, pp. 182--211, 1976. doi: 10.1147/sj.153.0182.

\bibitem{cohn2004} M. Cohn, \textit{User Stories Applied: For Agile Software
Development}. Boston, MA, USA: Addison-Wesley, 2004.

\bibitem{smart2014} J. F. Smart, \textit{BDD in Action: Behavior-Driven
Development for the Whole Software Lifecycle}. Shelter Island, NY, USA:
Manning Publications, 2014.

\bibitem{gotel1994} O. C. Z. Gotel and A. C. W. Finkelstein, ``An analysis of
the requirements traceability problem,'' in \textit{Proc. IEEE Int. Conf.
Requirements Engineering (ICRE)}, Colorado Springs, CO, USA, 1994,
pp. 94--101. doi: 10.1109/ICRE.1994.292398.

\bibitem{clelandhuang2012} J. Cleland-Huang, O. Gotel, and A. Zisman (Eds.),
\textit{Software and Systems Traceability}. London, U.K.: Springer, 2012.
doi: 10.1007/978-1-4471-2239-5.

\bibitem{wiegers2013} K. Wiegers and J. Beatty, \textit{Software
Requirements}, 3rd ed. Redmond, WA, USA: Microsoft Press, 2013.

% [COMPLETAR: agregar aquí las referencias adicionales que se citen en 3.1-3.5]

\end{thebibliography}

% ================================================================
% 14. ANEXOS
% ================================================================
\seccion{14.}{ANEXOS}

\subseccion{Anexo A -- Lista de verificación de inspección}
% Criterios ISO/IEC/IEEE 29148:2018. Una copia por inspector, firmada y fechada.

\textbf{Inspector:} Herrera Ramos, Thais Melanie \hfill \textbf{Rol:} Inspector 1

\begin{longtable}{p{2.0cm}p{3.2cm}p{1.6cm}p{1.0cm}p{4.0cm}}
\toprule
\textbf{Propiedad} & \textbf{Pregunta de verificación} & \textbf{Requisito} & \textbf{Cumple} & \textbf{Evidencia / Observación} \\
\midrule
\endhead
Trazabilidad & ¿Todo RNF tiene un RC (origen) asociado hacia atrás? & RNF-11--15 & N & Figuran como ``Elaboración propia'', sin RC asociado. \\
Factibilidad & ¿Se documentan los riesgos de dependencia de terceros? & RF-10/12 & N & Dependen de WhatsApp Business API; riesgo no documentado. \\
Consistencia & ¿La métrica del RNF coincide con lo exigido por su RF asociado? & RF-01/RNF-08 & N & RF-01 exige 100\,\%; RNF-08 exige $\geq$95\,\%. \\
Ambigüedad & ¿El umbral cuantificado aparece en el propio RF o solo en el RNF? & RF-05 & N & Remite implícitamente a RNF-03 sin cuantificar en el propio texto. \\
Trazabilidad & ¿El RF cita explícitamente al RNF que amplía su criterio? & RF-17/RNF-09 & N & No cita a RNF-09, rompiendo el patrón de RF-05/08/22. \\
Factibilidad & ¿La carga de usuarios concurrentes está respaldada por un RF? & RNF-11 & N & ``20 usuarios concurrentes'' sin RF que dimensione el personal. \\
\bottomrule
\end{longtable}
\textbf{Firma:} \rule{4cm}{0.4pt} \hfill \textbf{Fecha:} 05/08/2026

\vspace{0.4cm}
\textbf{Inspector:} Arboleda Yanza, Francisco Javier \hfill \textbf{Rol:} Lector

\begin{longtable}{p{2.0cm}p{3.2cm}p{1.6cm}p{1.0cm}p{4.0cm}}
\toprule
\textbf{Propiedad} & \textbf{Pregunta de verificación} & \textbf{Requisito} & \textbf{Cumple} & \textbf{Evidencia / Observación} \\
\midrule
\endhead
Consistencia & ¿El criterio de aceptación introduce funcionalidad no descrita en el RF? & RF-25 & N & Menciona selector de período no descrito en el RF. \\
Consistencia & ¿Los RF con lógica similar aplican el mismo control de validación? & RF-14 vs 17/19 & N & Sugerencias de IA exigen aprobación; recomendaciones nutricionales no. \\
Ambigüedad & ¿Se indica cuáles campos del formulario son obligatorios? & RF-04 & N & No aclara si los 10 campos son todos obligatorios. \\
Consistencia & ¿Se justifica por qué un canal se restringe respecto a otro RF similar? & RF-09 vs 10 & N & RF-10 restringe a WhatsApp sin justificar la diferencia. \\
Consistencia & ¿Dos RF que regulan el mismo mecanismo están diferenciados? & RF-22 vs 24 & N & No queda claro si son un mismo requisito fragmentado. \\
Trazabilidad & ¿El RF cita al RNF que amplía su criterio con una métrica adicional? & RF-08/RNF-02 & N & RF-08 no cita a RNF-02 en su descripción. \\
\bottomrule
\end{longtable}
\textbf{Firma:} \rule{4cm}{0.4pt} \hfill \textbf{Fecha:} 05/08/2026

\vspace{0.4cm}
\textbf{Inspector:} Marcillo Ponce, Alberto Jeanpool \hfill \textbf{Rol:} Moderador / Inspector 2

\begin{longtable}{p{2.0cm}p{3.2cm}p{1.6cm}p{1.0cm}p{4.0cm}}
\toprule
\textbf{Propiedad} & \textbf{Pregunta de verificación} & \textbf{Requisito} & \textbf{Cumple} & \textbf{Evidencia / Observación} \\
\midrule
\endhead
Verificabilidad & ¿El criterio de aceptación es medible con datos/umbrales concretos? & RF-16 & N & ``< 5 segundos'' no define tamaño de historial de referencia. \\
Completitud & ¿El requisito cubre los casos límite o excepcionales relevantes? & RF-21 & N & No especifica el comportamiento de RF-17--19 sin conexión. \\
Verificabilidad & ¿Se especifica quién ejecuta la acción y bajo qué rango? & RF-06 & N & ``Umbral mínimo configurable'' sin responsable ni rango. \\
Completitud & ¿Se define el flujo de excepción cuando la acción esperada no ocurre? & RF-15 & N & No indica qué pasa si el recordatorio no se atiende. \\
Verificabilidad & ¿Se define responsable y periodicidad de verificación de cumplimiento? & RNF-12 & N & No define quién audita el cifrado TLS/AES-256 ni con qué periodicidad. \\
Completitud & ¿El requisito contempla el caso de datos insuficientes (primera vez)? & RF-13/14 & N & RF-14 no define el caso sin historial suficiente. \\
\bottomrule
\end{longtable}
\textbf{Firma:} \rule{4cm}{0.4pt} \hfill \textbf{Fecha:} 05/08/2026

\subseccion{Anexo B -- Registro de defectos}
Ver la tabla consolidada de la sección 4.4 (18 defectos, D-01 a D-18) y la
tabla de correcciones aplicadas de la sección 5, ambas con el detalle
completo exigido por esta plantilla (N.\textordmasculine, sección, descripción,
tipo, severidad, requisito afectado y corrección).

\subseccion{Anexo C -- Formulario RFC}
Las tres instancias completas de RFC-01, RFC-02 y RFC-03, junto con el Acta
del CCB, están transcritas íntegramente en la sección 6 de este informe
(6.1--6.4). Las firmas autógrafas del Acta quedan pendientes de digitalizar
por el equipo antes de la entrega final.

\subseccion{Anexo D -- Exportación CSV del backlog}
\textcolor{uteqgris}{[PENDIENTE: exportar \texttt{backlog\_export.csv} desde
Jira \emph{después} de corregir la numeración de Stories descrita en 7.1, e
incluirlo en \texttt{04\_Trazabilidad/} del repositorio.]}

\subseccion{Anexo E -- Capturas}
\textcolor{uteqgris}{[PENDIENTE: capturas del tablero CASE completo, detalle de
$\geq$3 issues con enlaces visibles, panel de estadísticas (ver sección 7.4), y
del repositorio Git -- \texttt{git log --oneline --graph --decorate} y
\texttt{git tag -n} (ver sección 8.5).]}
% Ejemplo de inserción de imagen:
% \begin{figure}[htbp]
%   \centering
%   \includegraphics[width=0.85\textwidth]{figuras/tablero_completo.png}
%   \caption{Tablero CASE completo -- proyecto SGCV-IA}
% \end{figure}

\subseccion{Anexo F -- Declaración de uso de IA}
\begin{longtable}{p{4.3cm}p{8.5cm}}
\toprule
\textbf{Campo} & \textbf{Contenido} \\
\midrule
\endhead
Herramienta y versión & Claude (Anthropic), asistente conversacional con
herramientas de creación de documentos. \\
Tarea asistida & Generación del esqueleto LaTeX institucional del informe;
redacción del borrador del fundamento teórico (3.1--3.5) a partir de las
fuentes primarias provistas por el docente en la guía; transcripción y
formateo en LaTeX del contenido ya redactado por el equipo en los Anexos A, B
y C (checklists de inspección, registro de defectos, RFC y Acta del CCB) y en
el Mapa de Stakeholders; redacción de las cinco conclusiones críticas de la
sección 11 a partir de los resultados propios documentados en las secciones
4--6. \\
Fragmento afectado & Secciones 1--2 (carátula y resumen), 4--6 (inspección,
correcciones, CCB), 11 (conclusiones) y 12 (distribución del trabajo); las
secciones 3, 9 y 10 (fundamento teórico, comparativa CASE e IR
tradicional/ágil) fueron generadas en una sesión previa del equipo con el
mismo asistente. \\
Verificación crítica realizada por el equipo & Todos los defectos, RFC,
decisiones del CCB, roles y fechas transcritos en este informe provienen de
los archivos originales elaborados por el equipo (\texttt{AnexoA\_*.docx},
\texttt{AnexoB\_Registro\_Defectos\_completo.xlsx}, \texttt{RFC-01/02/03.docx},
\texttt{ACTA DEL CHANGE CONTROL BOARD.docx}, \texttt{Mapa\_de\_stakeholders.csv}),
no fueron generados ni inventados por el asistente. El equipo debe, antes de
la entrega: (1) corregir la numeración de Stories del tablero Jira conforme a
lo señalado en la sección 7.1, (2) completar los apartados marcados como
``PENDIENTE'' en las secciones 7, 8, 12 y en este mismo Anexo F con evidencia
real (capturas, CSV, historial de Git, firmas), y (3) leer íntegramente el
informe resultante antes de firmarlo, conforme al criterio A4 de la guía. \\
Firmas de los integrantes & Marcillo Ponce, Alberto Jeanpool; Herrera Ramos,
Thais Melanie; Arboleda Yanza, Francisco Javier
\textit{(firmas pendientes)}. \\
\bottomrule
\end{longtable}

\end{document}
